%%=============================================================================
%% Voorwoord
%%=============================================================================

\chapter*{\IfLanguageName{dutch}{Woord vooraf}{Preface}}%
\label{ch:voorwoord}

%% TODO:
%% Het voorwoord is het enige deel van de bachelorproef waar je vanuit je
%% eigen standpunt (``ik-vorm'') mag schrijven. Je kan hier bv. motiveren
%% waarom jij het onderwerp wil bespreken.
%% Vergeet ook niet te bedanken wie je geholpen/gesteund/... heeft

Ik heb dit onderwerp gekozen omdat de online wereld een overvloed aan informatie biedt die verre van optimaal benut wordt. De uitdaging om ongestructureerde data van diverse bronnen te centraliseren en bruikbaar te maken voor eindgebruikers, interesseert mij zeer veel. 
Webscraping is een technologie die zich in het spanningsveld tussen technische complexiteit en ethische verantwoordelijkheid bevindt, wat het voor mij een boeiend onderzoeksdomein maakt. Het ontwikkelen van een robuust systeem dat bestand is tegen de moderne defensieve mechanismen van het web, bood de perfecte gelegenheid om zowel mijn architecturale als technische vaardigheden verder te ontwikkelen.

Bovendien spreekt het informeren van concertgangers in een versnipperde Vlaamse markt, mij persoonlijk aan. De meerwaarde voor de eindgebruiker is wat mij motiveert om een gecentraliseerd notificatiesysteem te ontwikkelen.

Graag wil ik mijn promotor, mevrouw Heidi Roobrouck, en mijn co-promotor, de heer Jan Claes bedanken voor hun constructieve feedback en snelle beschikbaarheid. Beide waren van groot belang bij het verfijnen van dit onderzoek en het tot stand komen van de finale tekst. 
Ook wil ik mijn co-promotor, de heer Jan Claes, in het bijzonder bedanken voor het aanbieden en samen vormgeven van dit uitdagende onderwerp.

Tot slot wil ik mijn familie en vrienden bedanken voor hun steun tijdens dit proces. En mijn medestudenten voor de waardevolle discussies en gedeelde ervaringen die hebben bijgedragen aan zowel mijn persoonlijke groei als mijn visie tijdens dit onderzoek.


